% ======================================================================
% Module 4 Report: Physics-Informed Neural Networks for SEIRD
% A 3-page standalone report matching the style of report/main.tex.
% Compile with:  pdflatex module4_report  (twice for references)
%                bibtex module4_report
%                pdflatex module4_report (twice more)
% ======================================================================
\documentclass[11pt,a4paper]{article}

\usepackage[margin=2.5cm]{geometry}
\usepackage{amsmath,amssymb,amsthm}
\usepackage{booktabs}
\usepackage{graphicx}
\usepackage{hyperref}
\usepackage{parskip}
\usepackage{microtype}
\usepackage[T1]{fontenc}
\usepackage{lmodern}

\hypersetup{colorlinks=true,linkcolor=black,citecolor=black,urlcolor=blue}

% ── Theorem-like environments ────────────────────────────────────────────
\theoremstyle{remark}
\newtheorem*{remark}{Remark}

% ── Title block ─────────────────────────────────────────────────────────
\title{\textbf{Module 4: Physics-Informed Neural Networks for SEIRD Epidemiology}\\[4pt]
       {\large Mathematical Modelling Group Project --- Sierra Leone Ebola 2014--15}}
\author{Member 4}
\date{November 2026}

\begin{document}
\maketitle
\thispagestyle{empty}

% ======================================================================
\section{Introduction and Motivation}
% ======================================================================

Surveillance systems for Ebola report two cumulative quantities at irregular
dates: confirmed cases~$C(t)$ and confirmed deaths~$D(t)$.  Every other
epidemiologically relevant quantity---the susceptible fraction~$S$, the
incubating population~$E$, the currently infectious~$I$, and most critically
the time-varying transmission rate~$\beta(t)$---is latent.  Classical
approaches handle this latency by assuming a parametric shape for $\beta(t)$
in advance (a step function \cite{rivers2014ebola}, an exponential decay
\cite{althaus2014estimating}, or piecewise-constant segments).  Each
assumption has a qualitative cost: a step can never anticipate a second wave;
an exponential can only decrease.

Physics-informed neural networks (PINNs) \cite{raissi2019physics} offer a
different strategy.  Two small neural networks jointly represent the state
trajectories and $\beta(t)$ continuously, and are trained to minimise a loss
function that combines fidelity to the observed reports with a residual of the
SEIRD ordinary differential equations evaluated by automatic differentiation.
No shape for $\beta(t)$ is assumed.  The ODE constraint propagates information
from the observations into the hidden compartments: if the network fits the
reported cases \emph{and} satisfies the ODE, it must also be tracking~$E$
and~$I$ correctly.

This report describes the PINN architecture (Section~\ref{sec:arch}), the
loss function and training protocol (Section~\ref{sec:loss}), and the
quantitative results of three synthetic experiments designed to probe the
method's capabilities and limits (Section~\ref{sec:results}).

% ======================================================================
\section{Architecture and Inductive Biases}
\label{sec:arch}
% ======================================================================

\subsection{Observable Re-parameterisation}

Let $C(t) = I(t) + R(t) + D(t)$ denote the cumulative number of symptomatic
individuals.  Because $\dot C = \sigma E$, the SEIRD system can be written
entirely in terms of $E$, $I$, $C$, $D$:
\begin{align}
    \dot E &= \beta(t)\tfrac{S}{N}I - \sigma E, &
    \dot I &= \sigma E - (\gamma+\mu)I, \label{eq:ei} \\
    \dot C &= \sigma E, &
    \dot D &= \mu I, \label{eq:cd}
\end{align}
with $R = C-I-D$ and $S = N-E-C$.  The identity $S+E+I+R+D=N$ now holds for
\emph{any} choice of $(E,I,C,D)$, so the conservation law is built into the
parameterisation and cannot be violated.

Time is rescaled to $\tau = (t-t_0)/T \in [0,1]$ and every state is measured
in units of $K$, the largest cumulative case count in the training window
($K=1$ for synthetic fractions).  The equations become $\dot x = T f(x,\tau)$
with all quantities of order one, which stabilises the gradient of the
physics loss.

\subsection{Dual-Network Architecture}

\textbf{State network} $g_\theta(\tau)$: three hidden layers of 48 hyperbolic
tangent units, with four scalar outputs.  The states are:
\begin{equation}
    e(\tau) = e_0\,e^{\tau g_1(\tau)}, \quad
    i(\tau) = i_0\,e^{\tau g_2(\tau)}, \quad
    c(\tau) = c_0\,e^{\tau g_3(\tau)}, \quad
    d(\tau) = d_0 + \tau\, g_4(\tau).
    \label{eq:states}
\end{equation}
The exponential form ensures $e,i,c > 0$ for all $\tau$, and since each
function equals its initial value at $\tau=0$ exactly, the initial conditions
are satisfied without any penalty term.

\textbf{Transmission network} $h_\phi(\tau)$: three hidden layers of 32
hyperbolic tangent units, with $\beta(\tau) = \operatorname{Softplus}(h_\phi(\tau)+b) > 0$.

\subsection{Inductive Biases}

Six structural constraints are built into the model prior to training.
\begin{enumerate}
    \item \textbf{Hard initial conditions} (equation~\eqref{eq:states}).
    \item \textbf{Positivity of all states} (exponential and Softplus forms).
    \item \textbf{CFR bounds}: $\mu = \gamma\,\operatorname{CFR}/(1-\operatorname{CFR})$ with
          $\operatorname{CFR}\in(0.2,0.8)$, blocking the degenerate solution $\mu\to 0$.
    \item \textbf{Fourier features}: both networks receive $\tau' = 2\tau-1$
          together with $\sin(k\pi\tau')$ and $\cos(k\pi\tau')$ for
          $k=1,\ldots,4$.  Plain $\tanh$ networks preferentially learn
          low-frequency functions first (spectral bias \cite{rahaman2019spectral});
          the Fourier inputs allow rapid changes in $\beta(t)$ to be
          represented at any timescale within the window.
    \item \textbf{Relative ODE residuals}: each equation is divided by its
          state before squaring, so the early (small) phase of the epidemic
          is weighted equally with the late (large) phase.
    \item \textbf{Smooth prior on $\beta$}: a penalty $(d\beta/d\tau)^2$
          regularises the inverse problem, which would otherwise amplify
          measurement noise into spurious oscillations in $\beta$.
\end{enumerate}
Parameters $\sigma = 1/9.4\,\text{day}^{-1}$ and $\gamma = 1/7\,\text{day}^{-1}$
are fixed from clinical literature and are not estimated; Section~4.6 of the
main report explains the structural identifiability argument.

% ======================================================================
\section{Loss Function and Training Protocol}
\label{sec:loss}
% ======================================================================

\subsection{Loss Function}

The total loss is
\begin{equation}
    \mathcal{L} = \mathcal{L}_C + \mathcal{L}_D + \mathcal{L}_{\text{ode}} +
                  w_{\text{sm}}\,\mathcal{L}_{\text{smooth}},
    \label{eq:loss}
\end{equation}
where
\begin{align}
    \mathcal{L}_C &= \frac{1}{n_C}\sum_{k=1}^{n_C}\bigl(c(\tau_k) - C_k/K\bigr)^2,
    \qquad
    \mathcal{L}_D = \frac{1}{n_D}\sum_{k=1}^{n_D}\bigl(d(\tau_k) - D_k/K\bigr)^2,
    \label{eq:ldata}\\
    \mathcal{L}_{\text{ode}} &= \frac{1}{N_f}\sum_{j=1}^{N_f}
        \left[\left(\frac{e' - Tf_e}{e}\right)^{\!2}
             +\left(\frac{i' - Tf_i}{i}\right)^{\!2}
             +\left(\frac{c' - Tf_c}{c}\right)^{\!2}
             +\left(\frac{d' - Tf_d}{d}\right)^{\!2}\right]_{\!\tau_j},
    \label{eq:lode}\\
    \mathcal{L}_{\text{smooth}} &= \frac{1}{N_f}\sum_{j=1}^{N_f}
        \left(\frac{d\beta}{d\tau}(\tau_j)\right)^{\!2}.
    \label{eq:lsm}
\end{align}
All derivatives in~\eqref{eq:lode}--\eqref{eq:lsm} are computed by automatic
differentiation at $N_f = 200$ collocation points.  The sums in
$\mathcal{L}_C$ and $\mathcal{L}_D$ run only over real report dates; no
interpolation is used.  Weights are $w_{\text{sm}} = 10^{-4}$ for the
synthetic experiments and $10^{-3}$ for real-data fits.

\begin{remark}[Why relative residuals?]
An Ebola epidemic grows by two orders of magnitude over its window.  With
absolute residuals the early phase, where $e,i,c$ are small, contributes
almost zero loss.  The E6 ablation (Table~\ref{tab:e6}) shows that replacing
relative with absolute residuals raises the forward-problem errors from
below $2\times10^{-3}$ to 75--128\%.
\end{remark}

\subsection{Two-Stage Training}

Training proceeds in two stages.
\begin{itemize}
    \item \textbf{Adam} \cite{kingma2015adam} with learning rate $2\times10^{-3}$
          (cosine decay), $2{,}000$ iterations, fresh random collocation points
          each iteration.  This fast stochastic phase escapes saddle points and
          brings the loss into a basin of convergence.
    \item \textbf{L-BFGS} \cite{liu1989limited} with a strong-Wolfe line search,
          $2{,}500$--$6{,}000$ iterations on a fixed linspace grid.  The
          quasi-Newton update provides near-second-order convergence rates and is
          essential for the inverse problem (Section~\ref{sec:ablation}).
\end{itemize}
All computations use double precision (float64).  Every experiment is repeated
with 5 independent random seeds and the median is reported as the point
estimate.

Forecasting uses the network states at the end of the training window, holds
$\beta$ at its mean over the last 14 training days, and integrates the SEIRD
equations forward with the DOP853 solver.  This assumption---that $\beta$
remains constant---is the principal limit of the forecast horizon.

% ======================================================================
\section{Synthetic Experiments and Results}
\label{sec:results}
% ======================================================================

\subsection{Experiment E2: Forward Problem}

With $\beta = 0.4858\,\text{day}^{-1}$ ($R_0=2$), $\mu=0.1\,\text{day}^{-1}$,
and initial conditions $(s,e,i,r,d)(0) = (0.99,0.005,0.005,0,0)$, the PINN
is trained on $\mathcal{L}_{\text{ode}}$ alone (no data term).  The mean
relative $L_2$ errors against the DOP853 reference over 3 seeds are
$3.7\times10^{-4}$ ($s$), $5.4\times10^{-4}$ ($e$), $5.1\times10^{-4}$
($i$), $1.1\times10^{-3}$ ($r$) and $7.9\times10^{-4}$ ($d$); the worst
compartment in any seed is $2.0\times10^{-3}$.  Conservation holds to
$2\times10^{-16}$ (machine zero).

\subsection{Experiments E3 and E4: Inverse Problem}

An intervention on day~30 cuts $\beta$ from $0.55$ to $0.25\,\text{day}^{-1}$
over a 120-day window.  The PINN sees only daily cumulative cases and deaths
corrupted by independent multiplicative Gaussian noise of 1\% or 5\%.

\begin{table}[h]
\centering
\caption{E3/E4: recovery from noisy cumulative reports (mean $\pm$ SD, 5 seeds).
Errors are relative $L_2$ except $\mu$ and RMSE rows.}
\label{tab:e3e4}
\begin{tabular}{lcc}
\toprule
Quantity & 1\% noise & 5\% noise \\
\midrule
$\beta(t)$ rel.\ $L_2$                            & $11.0\%\pm0.7\%$   & $11.6\%\pm0.6\%$   \\
$\beta(t)$ RMSE away from the jump (day$^{-1}$)   & $0.022\pm0.005$    & $0.025\pm0.005$    \\
$\mu$ relative error                              & $4.6\%\pm0.3\%$    & $5.4\%\pm1.0\%$    \\
hidden $s(t)$                                     & $0.36\%\pm0.03\%$  & $0.46\%\pm0.09\%$  \\
hidden $e(t)$                                     & $5.8\%\pm0.6\%$    & $6.4\%\pm0.7\%$    \\
hidden $i(t)$                                     & $5.8\%\pm0.4\%$    & $6.5\%\pm0.8\%$    \\
\bottomrule
\end{tabular}
\end{table}

The transmission rate is recovered within 11--12\% overall; most of that
error is concentrated at the abrupt jump, which a smooth continuous network
cannot follow instantaneously.  The death rate $\mu$ is recovered within 5\%
despite being unobserved directly: the ratio of the death curve slope to the
case curve slope carries its value.  The susceptible fraction is recovered
within 0.5\%, and the latent compartments $E$ and $I$ within 6.5\%.
Multiplying the noise level fivefold costs less than one percentage point on
every metric, demonstrating robustness.

\subsection{Ablation Study (Experiment E6)}
\label{sec:ablation}

\begin{table}[h]
\centering
\caption{E6: ablation on the E3 problem (5\% noise, mean of 2 seeds).}
\label{tab:e6}
\begin{tabular}{lccc}
\toprule
Setting & $\beta(t)$ RMSE & hidden $e(t)$ rel.\ $L_2$ & $\mu$ rel.\ error \\
\midrule
\textbf{Full method}                            & \textbf{0.042} & \textbf{0.070} & 0.054 \\
$w_{\text{ode}} = 10$ (instead of 1)           & 0.123          & 0.207          & 0.143 \\
Adam only (8000 iters, no L-BFGS)              & 0.243          & 0.957          & 0.176 \\
No Fourier features                             & 0.058          & 0.083          & 0.045 \\
\bottomrule
\end{tabular}
\end{table}

\textbf{The L-BFGS stage is essential.} Without it, the exposed compartment
error is 96\% and $\beta(t)$ stays nearly flat.  \textbf{Loss balance
matters}: increasing the physics weight tenfold pushes the network to satisfy
the ODE at the expense of the data, roughly tripling every error.
\textbf{Fourier features} reduce the $\beta$ RMSE by approximately 28\%.
\textbf{Relative ODE residuals} are decisive for the forward problem (E2):
switching to absolute residuals gives errors of 75--128\% instead of below
$2\times10^{-3}$.

\section*{Conclusion}

The PINN demonstrates that a smooth nonparametric $\beta(t)$, unobserved
compartments ($E$, $I$, $S$), and the death rate $\mu$ can all be inferred
reliably from cumulative case and death reports alone, provided the ODE
constraint is enforced by autograd-based collocation.  The six inductive
biases (hard initial conditions, positivity, CFR bounds, Fourier features,
relative residuals, and $\beta$-smoothness) together explain the gap between
the full method and its ablated variants.  The principal scientific limit
is the forecasting assumption that $\beta$ remains constant beyond the
training window.

% ── Bibliography ─────────────────────────────────────────────────────────
\begin{thebibliography}{9}
\bibitem{raissi2019physics}
  M.~Raissi, P.~Perdikaris, and G.~E.~Karniadakis,
  ``Physics-informed neural networks: A deep learning framework for solving forward and inverse problems involving nonlinear partial differential equations,''
  \emph{J.\ Comput.\ Phys.}, 378:686--707, 2019.

\bibitem{rivers2014ebola}
  C.~Rivers \emph{et al.},
  ``Modeling the impact of interventions on an epidemic of Ebola in Sierra Leone and Liberia,''
  \emph{PLoS Curr.\ Outbreaks}, 2014.

\bibitem{althaus2014estimating}
  C.~L.~Althaus,
  ``Estimating the reproduction number of Ebola virus (EBOV) during the 2014 outbreak in West Africa,''
  \emph{PLoS Curr.\ Outbreaks}, 2014.

\bibitem{rahaman2019spectral}
  N.~Rahaman \emph{et al.},
  ``On the spectral bias of neural networks,''
  \emph{ICML}, 2019.

\bibitem{kingma2015adam}
  D.~P.~Kingma and J.~Ba,
  ``Adam: A method for stochastic optimization,''
  \emph{ICLR}, 2015.

\bibitem{liu1989limited}
  D.~C.~Liu and J.~Nocedal,
  ``On the limited memory BFGS method for large scale optimization,''
  \emph{Math.\ Programming}, 45(1):503--528, 1989.
\end{thebibliography}

\end{document}
